\documentclass[12pt]{article}
\usepackage[T1]{fontenc}
\usepackage[utf8]{inputenc}
\usepackage{lmodern}
\usepackage{textcomp}
\usepackage{enumitem}
\usepackage{geometry}
\usepackage{graphicx}
\usepackage{amsmath, amssymb}
\geometry{margin=1in}
\begin{document}
\begin{center}
Chemistry - Graduate Level Assessment 8 \
Total Marks: 40 \
Answer all questions.
\end{center}

\section*{Multiple Choice (10 marks)}
\begin{enumerate}[label=\arabic*.]
\item What pressure would $131 \mathrm{g}$ of xenon gas in a vessel of volume $1.0 \mathrm{dm}^3$ exert at $25^{\circ} \mathrm{C}$ if it behaved as a van der Waals gas?
\begin{enumerate}[label=(\alph*)]
    \item 22$\mathrm{atm}$
    \item 20$\mathrm{atm}$
    \item 15$\mathrm{atm}$
    \item 30$\mathrm{atm}$
    \item 34$\mathrm{atm}$
\end{enumerate}
\item Determine the equilibrium constant for the reaction H\_2 +I\_2 \ensuremath{\rightleftarrows} 2 HI the equilibrium concentrations are: H\_2 , 0.9 moles/liter; I\_2 , 0.4 mole/liter; HI, 0.6 mole/liter.
\begin{enumerate}[label=(\alph*)]
    \item 0.9
    \item 1.0
    \item 2.5
    \item 0.3
    \item 0.5
\end{enumerate}
\item The temperature of the fireball in a thermonuclear explosion can reach temperatures of approximately $10^7 \mathrm{~K}$. What value of $\lambda_{\max }$ does this correspond to?
\begin{enumerate}[label=(\alph*)]
    \item 2 $10^{-10} \mathrm{~m}$
    \item 1 $10^{-10} \mathrm{~m}$
    \item 7 $10^{-10} \mathrm{~m}$
    \item 6 $10^{-10} \mathrm{~m}$
    \item 4 $10^{-10} \mathrm{~m}$
\end{enumerate}
\item The density of a gaseous compound was found to be $1.23 \mathrm{kg} \mathrm{m}^{-3}$ at $330 \mathrm{K}$ and $20 \mathrm{kPa}$. What is the molar mass of the compound?
\begin{enumerate}[label=(\alph*)]
    \item 169 $\mathrm{g} \mathrm{mol}^{-1}$
    \item 123 $\mathrm{g} \mathrm{mol}^{-1}$
    \item 200 $\mathrm{g} \mathrm{mol}^{-1}$
    \item 150 $\mathrm{g} \mathrm{mol}^{-1}$
    \item 190 $\mathrm{g} \mathrm{mol}^{-1}$
\end{enumerate}
\item A liquid element that is a dark-colored, nonconducting substance at room temperature is
\begin{enumerate}[label=(\alph*)]
    \item lead
    \item cesium
    \item bromine
    \item francium
    \item rubidium
\end{enumerate}
\end{enumerate}

\section*{True / False (10 marks)}
\textit{Answer True or False}
\begin{enumerate}[label=\arabic*.]
\setcounter{enumi}{5}
\item True or False: The height of a mercury column corresponding to a 600 mm water column is approximately 44.1 mm Hg due to the density difference.
\item True or False: The positron is the most massive subatomic particle among the listed options.
\item True or False: The logarithm of 0.0364 is approximately -1.4389.
\item True or False: An electron transitioning from energy level 4 to 1 emits a photon with a wavelength of $1.03 \times 10^{-7} \mathrm{~m}$.
\item True or False: The combustion of 1000 cubic feet of butane gas requires 5500 cubic feet of oxygen for complete reaction.
\end{enumerate}

\section*{Long Form Response (20 marks)}
\textit{Answer each question with a clear and complete explanation.}
\begin{enumerate}[label=\arabic*.]
\setcounter{enumi}{10}
\item Calculate the volume of concentrated sulfuric acid needed to achieve a mass of 74.0 g, given its density of 1.85 g/ml. Discuss the implications of density in this context.
\item A chemist dropped a 200 g object into water, displacing 60 ml, and later displaced 50 g of oil. Calculate and analyze the densities of both the object and the oil, discussing the implications of your findings.
\end{enumerate}

\vfill
\noindent\textit{End of Paper}
\end{document}